\subsection{Corrigés des exercices}

\targetCorr{fsec2_ex1}
\begin{correctionbox}[Corrigé de l'exercice 1 : Arbitrage Parité Call-Put]
Calcul de la relation théorique : $C - P = 10 - 8 = 2$.
Calcul de $S_0 - K e^{-rT} = 100 - 100 \times e^{-0,05} \approx 100 - 95,12 = 4,88$.
Puisque $2 < 4,88$, la relation n'est pas respectée. Le côté gauche (les options) est sous-évalué par rapport au côté droit.

Portefeuille d'arbitrage à t=0 :
On achète ce qui est moins cher ($C - P$) et on vend ce qui est plus cher ($S_0 - K e^{-rT}$) :
\begin{itemize}[label=$\cdot$]
    \item Acheter le Call (-10 €) et vendre le Put (+8 €).
    \item Vendre l'action à découvert (+100 €) et prêter la différence au taux sans risque.
    \item Trésorerie nette placée : $100 + 8 - 10 = 98$ €.
\end{itemize}
À T=1 :
Le placement rapporte $98 \times e^{0,05} \approx 103,02$ €.
Si $S_T > 100$, on exerce le Call : on achète l'action à 100 pour la rendre. Coût : -100 €.
Si $S_T \le 100$, l'acheteur exerce le Put : on doit lui acheter l'action à 100. Coût : -100 €.
Dans les deux cas, le coût final est de 100 €.
Profit net : $103,02 - 100 = 3,02$ € (sans risque).
\navToExo{fsec2_ex1}
\end{correctionbox}

\targetCorr{fsec2_ex2}
\begin{correctionbox}[Corrigé de l'exercice 2 : Borne inférieure du Call]
Considérons deux portefeuilles :
Portefeuille A : 1 Call européen + $K e^{-rT}$ investis au taux sans risque.
Valeur à l'échéance $T$ : $\max(S_T - K, 0) + K = \max(S_T, K)$.
Portefeuille B : 1 Action $S$.
Valeur à l'échéance $T$ : $S_T$.

Puisque $\max(S_T, K) \geq S_T$ pour tout $S_T$, le Portefeuille A vaudra toujours au moins autant que le Portefeuille B à l'échéance. Par absence d'arbitrage, cette inégalité doit tenir à $t=0$ :
$C_0 + K e^{-rT} \geq S_0 \implies C_0 \geq S_0 - K e^{-rT}$.
Comme un Call ne peut avoir un prix négatif, on a bien $C_0 \geq \max(0, S_0 - K e^{-rT})$.
\navToExo{fsec2_ex2}
\end{correctionbox}

\targetCorr{fsec2_ex3}
\begin{correctionbox}[Corrigé de l'exercice 3 : Borne inférieure du Put]
Par la parité Call-Put : $P_0 = C_0 - S_0 + K e^{-rT}$.
Puisque le prix d'un Call est toujours positif ou nul ($C_0 \ge 0$), nous avons :
$P_0 \ge 0 - S_0 + K e^{-rT} \implies P_0 \ge K e^{-rT} - S_0$.
Comme une option a une responsabilité limitée, son prix est toujours non négatif. 
On obtient : $P_0 \ge \max(0, K e^{-rT} - S_0)$.
\navToExo{fsec2_ex3}
\end{correctionbox}

\targetCorr{fsec2_ex4}
\begin{correctionbox}[Corrigé de l'exercice 4 : Inégalité des écarts de Strike]
Construisons un Bull Spread : achetons le Call $K_1$ et vendons le Call $K_2$.
Le coût initial est $C_0(K_1) - C_0(K_2)$.
À l'échéance, le payoff de ce portefeuille est :
\begin{itemize}[label=$\cdot$]
    \item Si $S_T \le K_1$ : $0 - 0 = 0$.
    \item Si $K_1 < S_T \le K_2$ : $S_T - K_1 - 0 = S_T - K_1$ (le maximum est $K_2 - K_1$).
    \item Si $S_T > K_2$ : $(S_T - K_1) - (S_T - K_2) = K_2 - K_1$.
\end{itemize}
Le gain maximal de ce portefeuille à l'échéance est plafonné à $(K_2 - K_1)$.
Par absence d'arbitrage, le prix de ce portefeuille aujourd'hui ne peut pas excéder la valeur actualisée de son gain maximal garanti :
$C_0(K_1) - C_0(K_2) \le (K_2 - K_1)e^{-rT}$.
\navToExo{fsec2_ex4}
\end{correctionbox}